%% toc: P2 23 23.5,23.6,23.7,23.8
\section{Proofs for the direct nonenumerability of full second-order validity}
\label{sec:annex-p02-w02-p02-041}

Notation is that of
\cref{def:p02-w02-p02-041-terms,def:p02-w02-p02-041-graph-sentence}.
\(\mathcal N=(\N,0,n\mapsto n+1,+,\times)\) is the standard structure, and
\(\overline n\) is the numeral \(S\cdots S0\) with \(n\) occurrences of \(S\).
Induction on \(n\) shows that \(\overline n\) denotes \(n\) in \(\mathcal N\)
under every assignment.

\begin{annexitem}{annex:p02-w02-p02-041-coverage}{thm:p02-w02-p02-041-coverage}{Every recursive function has a term}
\begin{proof}
Let \(\mathcal F\) be the class of partial maps \(f_\rho\), for \(\rho\)
ranging over terms, each taken with arity \(\mathrm{ar}(\rho)\).

\emph{Initial functions.}  \(\mathcal F\) contains \(Z_k=f_{\mathbf Z^k}\) for
every \(k\in\N\), the successor \(f_{\mathbf S}\), and every projection
\(P^k_i=f_{\mathbf P^k_i}\) with \(1\le i\le k\).

\emph{Closure.}  Let \(g=f_\gamma\) have arity \(m\) and
\(h_1=f_{\eta_1},\ldots,h_m=f_{\eta_m}\) have arity \(k\).  Then
\(\mathbf C^k(\gamma;\eta_1,\ldots,\eta_m)\) is a term (its arity side
conditions are exactly these arities), and by
\cref{def:p02-w02-p02-041-terms} it denotes the composition with the
definedness convention of \cref{def:p04-partial-recursive}.  If
\(g=f_\gamma\) has arity \(k\) and \(h=f_\eta\) arity \(k+2\), the term
\(\mathbf R(\gamma,\eta)\) denotes the map defined from \(g,h\) by primitive
recursion, with the same definedness convention, for every \(k\in\N\).  If
\(g=f_\gamma\) has arity \(k+1\), \(\mathbf M(\gamma)\) denotes
\(\mu g\) as defined in \cref{def:p04-partial-recursive}.

(a) The members of \(\mathcal F\) of positive arity include the zero,
successor and projection functions of positive arity and are closed under
composition, primitive recursion and unbounded minimization whenever the
result has positive arity, because each operation applied to members of
\(\mathcal F\) yields a member of \(\mathcal F\) of the same arity.  The
partial recursive functions are the least class with these properties, so
each of them is in \(\mathcal F\) with its arity.

(b) If \(g,h_1,\ldots,h_m\) are total, the composition is total and is the
composition of \cref{sec:turing-machines-02}; if \(g,h\) are total, induction
on \(n\) shows that the map defined by primitive recursion is total, and it
satisfies the two equations of \cref{sec:turing-machines-02}.  So the total
members of \(\mathcal F\) contain \(Z_k\) for every \(k\in\N\) (including
\(Z_0\)), the successor and the projections, and are closed under the two
operations of \cref{sec:turing-machines-02}, for every \(k\in\N\).  The
primitive recursive functions are the least such class, so each is in
\(\mathcal F\) with its arity.
\end{proof}
\end{annexitem}

\begin{annexitem}{annex:p02-w02-p02-041-iso}{thm:p02-w02-p02-041-iso}{Isomorphism invariance under full semantics}
\begin{proof}
\emph{Terms.}  For every term \(t\),
\(f(\sem{t}^{\mathcal A}_s)=\sem{t}^{\mathcal B}_{f\circ s}\), by induction on
\(t\): for a variable this is the definition of \(f\circ s\); for a constant
and for a function symbol applied to terms it is the preservation of constants
and functions by \(f\) together with the induction hypothesis.

\emph{Formulas,} by induction on \(\varphi\), simultaneously for all
assignments \(s\).  For \(t_1=t_2\): \(f\) is injective, so
\(\sem{t_1}^{\mathcal A}_s=\sem{t_2}^{\mathcal A}_s\) iff their images are
equal, which by the term clause are the values in \(\mathcal B\).  For a
relation symbol \(P\) of the signature, \(f\) preserves and reflects \(P\).
For a relation-variable atom \(X(t_1,\ldots,t_n)\) put
\(a_i=\sem{t_i}^{\mathcal A}_s\).  The values in \(\mathcal B\) are
\(f(a_1),\ldots,f(a_n)\).  If \((a_1,\ldots,a_n)\in s(X)\) then
\((f(a_1),\ldots,f(a_n))\in(f\circ s)(X)\) by definition.  Conversely, if
\((f(a_1),\ldots,f(a_n))=(f(b_1),\ldots,f(b_n))\) for some
\((b_1,\ldots,b_n)\in s(X)\), injectivity gives \(a_i=b_i\) for each \(i\), so
\((a_1,\ldots,a_n)\in s(X)\).  The constant \(\bottom\) and the connectives
follow from the induction hypothesis.

For \(\forall x\,\psi\): \(f\circ(s[x\mapsto a])=(f\circ s)[x\mapsto f(a)]\), and
as \(a\) ranges over \(A\), \(f(a)\) ranges over all of \(B\) because \(f\) is
onto; so \(\psi\) holds for every \(a\) on the left iff it holds for every
\(b\) on the right.  The clause for \(\exists x\) is the same with ``some''.

For \(\forall X^n\psi\): write \(f[R]=\{(f(a_1),\ldots,f(a_n)):(a_1,\ldots,a_n)\in R\}\)
for \(R\subseteq A^n\).  Then
\(f\circ(s[X\mapsto R])=(f\circ s)[X\mapsto f[R]]\).  The map
\(R\mapsto f[R]\) is a bijection from \(\mathcal P(A^n)\) onto \(\mathcal P(B^n)\): its
inverse is \(R'\mapsto\{(a_1,\ldots,a_n):(f(a_1),\ldots,f(a_n))\in R'\}\),
because \(f\) is a bijection.  Under full semantics the quantifier on the left
ranges over all of \(\mathcal P(A^n)\) and on the right over all of \(\mathcal P(B^n)\),
so the induction hypothesis for \(\psi\) gives the clause; \(\exists X^n\) is
the same.
\end{proof}
\end{annexitem}

\begin{annexitem}{annex:p02-w02-p02-041-graphs}{thm:p02-w02-p02-041-graphs}{The graph conditions have exactly one solution}
\begin{proof}
\emph{Part 1.}  Let \(\sigma\) have arity \(k\), assume
\(s(X_\tau)=\mathrm{Gr}_\tau\) for its immediate subterms, and put
\(R=s(X_\sigma)\).  In \(\mathcal N\) the symbol \(0\) denotes \(0\), \(S\)
the successor and \(+\) addition, so the atomic parts of
\(\mathsf{Def}_\sigma\) have their arithmetic meaning.  In every case
\(\mathsf{Def}_\sigma\) has the form
\(\forall\vec x\,\forall x_0\,(X_\sigma(\vec x,x_0)\leftrightarrow\chi)\), or (for
\(\mathbf R\)) a conjunction of two such clauses, and we compute the condition
\(\chi\) expresses.

\(\mathbf Z^k\), \(\mathbf S\), \(\mathbf P^k_i\): the condition says
\((\vec x,y)\in R\) iff \(y=0\), respectively \(y=x_1+1\), respectively
\(y=x_i\), for all \(\vec x,y\); that is, \(R=\mathrm{Gr}_\sigma\).

\(\mathbf C^k(\gamma;\eta_1,\ldots,\eta_m)\): for given \(\vec x,y\), the right
side holds iff there are \(z_1,\ldots,z_m\) with
\(f_{\eta_j}(\vec x)\simeq z_j\) for each \(j\) and
\(f_\gamma(z_1,\ldots,z_m)\simeq y\).  Such \(z_j\) exist iff every
\(f_{\eta_j}(\vec x)\) is defined, and then they are those values; so the right
side holds iff \(f_\sigma(\vec x)\simeq y\).  Hence the condition says
\(R=\mathrm{Gr}_\sigma\).  (For \(m=0\) the right side is
\(X_\gamma(x_0)\) and \(f_\sigma(\vec x)\simeq f_\gamma()\).)

\(\mathbf M(\gamma)\), with \(\mathrm{ar}(\gamma)=k+1\): for \(z,y\in\N\),
\(\exists w\,(z+(w+1)=y)\) holds iff \(z<y\).  So the right side holds at
\(\vec x,y\) iff \(f_\gamma(\vec x,y)\simeq0\) and, for every \(z<y\), some
\(v\neq0\) has \(f_\gamma(\vec x,z)\simeq v\).  This is exactly the statement
that \(y\) is the value \(\mu y[f_\gamma(\vec x,y)=0]\) of
\cref{def:p02-w02-p02-041-terms}: such a \(y\) is the least zero with all
earlier values defined and nonzero, and conversely that value has the
property.  So the right side holds iff \(f_\sigma(\vec x)\simeq y\), and the
condition says \(R=\mathrm{Gr}_\sigma\).

\(\mathbf R(\gamma,\eta)\), with \(\mathrm{ar}(\gamma)=k\).  If
\(R=\mathrm{Gr}_\sigma\), the first clause holds because
\(f_\sigma(\vec x,0)\simeq f_\gamma(\vec x)\), and the second holds because
\(f_\sigma(\vec x,n+1)\simeq y\) iff some \(u\) has
\(f_\sigma(\vec x,n)\simeq u\) and \(f_\eta(\vec x,n,u)\simeq y\).  Conversely
assume both clauses.  We prove by induction on \(n\) that for all \(\vec x,y\),
\((\vec x,n,y)\in R\) iff \((\vec x,n,y)\in\mathrm{Gr}_\sigma\).  For
\(n=0\), the first clause gives \((\vec x,0,y)\in R\) iff
\(f_\gamma(\vec x)\simeq y\) iff \(f_\sigma(\vec x,0)\simeq y\).  For \(n+1\),
the second clause (with \(x_{k+1}\) valued \(n\), so that \(Sx_{k+1}\) has value
\(n+1\)) gives \((\vec x,n+1,y)\in R\) iff some \(u\) has
\((\vec x,n,u)\in R\) and \(f_\eta(\vec x,n,u)\simeq y\); by the induction
hypothesis \((\vec x,n,u)\in R\) iff \(f_\sigma(\vec x,n)\simeq u\); so the
condition is \(f_\sigma(\vec x,n+1)\simeq y\).  Every member of \(\N\) is
\(0\) or some \(n+1\), so \(R=\mathrm{Gr}_\sigma\).

\emph{Part 2.}  By \cref{def:p02-w02-p02-041-graph-sentence},
\(\mathrm{Sub}(\rho)\) contains the immediate subterms of each of its members.
If \(s(X_\sigma)=\mathrm{Gr}_\sigma\) for all \(\sigma\) in the list, part~1
applies to each member and gives \(\mathcal N,s\entails\mathsf{Def}_\sigma\).
Conversely assume every \(\mathsf{Def}_\sigma\) with \(\sigma\) in the list.
By induction on the code \(\ulcorner\sigma\urcorner\) for \(\sigma\) in the
list: the immediate subterms of \(\sigma\) are in the list and have smaller
codes (\cref{sec:arithmetic-coding-04}: list entries and constructor arguments
are less than the code), so by the induction hypothesis
\(s(X_\tau)=\mathrm{Gr}_\tau\) for each of them; part~1 then gives
\(s(X_\sigma)=\mathrm{Gr}_\sigma\).  Different subterms have different codes and
hence different variables \(X_\sigma\), so these requirements never conflict.
\end{proof}
\end{annexitem}

\begin{annexitem}{annex:p02-w02-p02-041-reduction}{thm:p02-w02-p02-041-reduction}{The divergence sentence is valid exactly at divergence}
\begin{proof}
Write \(c=\ulcorner\rho\urcorner\),
\(\mathrm{Sub}(\rho)=[\sigma_1,\ldots,\sigma_n]\) with \(\sigma_1=\rho\),
\(\Delta\) for \(\mathsf{Def}_{\sigma_1}\land\cdots\land\mathsf{Def}_{\sigma_n}\),
and \(\Omega\) for \(\neg\exists x_0\,X_\rho(\overline c,x_0)\), so that
\(\Theta_\rho\) is \(\mathsf{PA}_2\to\forall X_{\sigma_1}\cdots\forall X_{\sigma_n}(\Delta\to\Omega)\).

\emph{Validity implies divergence.}  Assume \(\entails_{\mathrm{full}}\Theta_\rho\).
By \cref{thm:p02-w02-p02-040-pa2-categorical}, \(\mathcal N\) satisfies the
sentence \(\mathsf{PA}_2\) under full semantics.  Take any full assignment
\(s\) over \(\mathcal N\) and let \(s'\) agree with \(s\) except that
\(s'(X_\sigma)=\mathrm{Gr}_\sigma\) for each \(\sigma\) in
\(\mathrm{Sub}(\rho)\); this is a full assignment, since full semantics allows
every relation.  Validity gives \(\mathcal N,s'\entails\Delta\to\Omega\), and
\cref{thm:p02-w02-p02-041-graphs}(2) gives \(\mathcal N,s'\entails\Delta\).
So \(\mathcal N,s'\entails\Omega\): there is no \(y\) with
\((c,y)\in\mathrm{Gr}_\rho\), since \(\overline c\) denotes \(c\).  Hence
\(f_\rho(c)\) is undefined.

\emph{Divergence implies validity.}  Assume \(f_\rho(c)\) is undefined.  Let
\(\mathcal M\) be any structure for \(\{0,S,+,\times\}\) and \(s\) any full
assignment over it.  If \(\mathcal M,s\) does not satisfy \(\mathsf{PA}_2\),
the implication \(\Theta_\rho\) holds.  Otherwise
\cref{thm:p02-w02-p02-040-pa2-categorical} gives an isomorphism \(f\) from
\(\mathcal N\) onto \(\mathcal M\), and \(f^{-1}\) is an isomorphism from
\(\mathcal M\) onto \(\mathcal N\).  Let \(s'\) differ from \(s\) at most at
\(X_{\sigma_1},\ldots,X_{\sigma_n}\), and suppose
\(\mathcal M,s'\entails\Delta\).  By \cref{thm:p02-w02-p02-041-iso} applied
to \(f^{-1}\), \(\mathcal N,f^{-1}\circ s'\entails\Delta\), so by
\cref{thm:p02-w02-p02-041-graphs}(2) the value of \(X_\rho\) under
\(f^{-1}\circ s'\) is \(\mathrm{Gr}_\rho\), which contains no pair with first
entry \(c\).  Hence \(\mathcal N,f^{-1}\circ s'\entails\Omega\), and
\cref{thm:p02-w02-p02-041-iso} again gives \(\mathcal M,s'\entails\Omega\).  As
\(s'\) was arbitrary, \(\mathcal M,s\) satisfies
\(\forall X_{\sigma_1}\cdots\forall X_{\sigma_n}(\Delta\to\Omega)\) (a repeated
variable in the prefix is quantified twice, which changes nothing), and so
satisfies \(\Theta_\rho\).  As \(\mathcal M\) and \(s\) were arbitrary,
\(\entails_{\mathrm{full}}\Theta_\rho\).
\end{proof}
\end{annexitem}

\begin{annexitem}{annex:p02-w02-p02-041-effective}{thm:p02-w02-p02-041-effective}{The divergence sentence is primitive recursive in the term}
\begin{proof}
All functions below are built from the primitive recursive list and
constructor operations of \cref{sec:arithmetic-coding-04}: \(C\), \(h\),
\(t\), length \(\ell\), entry \((c)_j\), concatenation \(\mathbin{+\!\!+}\),
truncated subtraction and the tests derived from it, the numeral code
\(\nu(n)=\ulcorner\overline n\urcorner\), the bounded-loop principle (a loop with
a coded-list state, a primitive recursive update and a primitive recursive
iteration bound is primitive recursive), and the course-of-values principle
(a function whose value at \(c\) is a primitive recursive combination of its
values at proper child codes of \(c\) is primitive recursive).  Definition by
cases on primitive recursive tests preserves primitive recursiveness, because a
case split is a sum of products with characteristic functions.

\emph{Step 1: term test and arity.}  By one course-of-values recursion define
\(\mathrm{RT}(c)\in\{0,1\}\) and \(\mathrm{ar}(c)\).  At row \(c\) put
\(k=h(c)\) and \(a=t(c)\), and for tag \(33\) put \(L=(a)_2\).
\begin{itemize}
 \item \(k=30\): \(\mathrm{RT}=1\) iff \(\ell(a)=1\); \(\mathrm{ar}=(a)_0\).
 \item \(k=31\): \(\mathrm{RT}=1\) iff \(\ell(a)=0\); \(\mathrm{ar}=1\).
 \item \(k=32\): \(\mathrm{RT}=1\) iff \(\ell(a)=2\) and
   \(1\le(a)_1\le(a)_0\); \(\mathrm{ar}=(a)_0\).
 \item \(k=33\): \(\mathrm{RT}=1\) iff \(\ell(a)=3\),
   \(\mathrm{RT}((a)_1)=1\), \(\mathrm{ar}((a)_1)=\ell(L)\), and for every
   \(j<\ell(L)\), \(\mathrm{RT}((L)_j)=1\) and \(\mathrm{ar}((L)_j)=(a)_0\);
   \(\mathrm{ar}=(a)_0\).
 \item \(k=34\): \(\mathrm{RT}=1\) iff \(\ell(a)=2\),
   \(\mathrm{RT}((a)_0)=\mathrm{RT}((a)_1)=1\) and
   \(\mathrm{ar}((a)_1)=\mathrm{ar}((a)_0)+2\); \(\mathrm{ar}=\mathrm{ar}((a)_0)+1\).
 \item \(k=35\): \(\mathrm{RT}=1\) iff \(\ell(a)=1\), \(\mathrm{RT}((a)_0)=1\)
   and \(\mathrm{ar}((a)_0)\ge1\); \(\mathrm{ar}=\mathrm{ar}((a)_0)\mathbin{dot{-}}1\).
 \item Any other row, including \(c=0\): \(\mathrm{RT}=0\), \(\mathrm{ar}=0\).
\end{itemize}
In every case put \(\mathrm{ar}=0\) whenever \(\mathrm{RT}=0\).  The rows consulted, \((a)_0\), \((a)_1\) and
\((L)_j\), are proper child codes: \(a=t(c)<c\), list entries are less than
the list, and \(L=(a)_2<a\).  The universal condition over \(j<\ell(L)\) is a
bounded loop.  So \(\mathrm{RT}\) and \(\mathrm{ar}\) are primitive recursive.
Induction on \(c\), using that \(c>0\) determines its tag and argument list
uniquely, shows that \(\mathrm{RT}(c)=1\) iff \(c\) is the code of a term, and
that \(\mathrm{ar}(c)\) is then its arity; the clauses are exactly those of
\cref{def:p02-w02-p02-041-terms}.

\emph{Step 2: subterm lists.}  By course-of-values recursion,
\(\mathrm{Sb}(c)=[c]\) at tags \(30\) to \(32\);
\(\mathrm{Sb}(c)=[c]\mathbin{+\!\!+}\mathrm{Sb}((a)_1)\mathbin{+\!\!+}\mathrm{Sb}((L)_0)\mathbin{+\!\!+}\cdots\mathbin{+\!\!+}\mathrm{Sb}((L)_{\ell(L)-1})\)
at tag \(33\), the iterated concatenation being a loop of length \(\ell(L)\);
\(\mathrm{Sb}(c)=[c]\mathbin{+\!\!+}\mathrm{Sb}((a)_0)\mathbin{+\!\!+}\mathrm{Sb}((a)_1)\) at
tag \(34\); \(\mathrm{Sb}(c)=[c]\mathbin{+\!\!+}\mathrm{Sb}((a)_0)\) at tag \(35\);
and \(\mathrm{Sb}(c)=[]\) at other rows.  For a term \(\rho\),
\(\mathrm{Sb}(\ulcorner\rho\urcorner)\) is the list of codes of
\(\mathrm{Sub}(\rho)\), by induction on \(\rho\).

\emph{Step 3: formula builders.}  The following are primitive recursive:
\(v(i)=N_0(i)\) (the variable \(x_i\));
\(V(p,0)=[]\), \(V(p,q+1)=V(p,q)\mathbin{+\!\!+}[v(p+q)]\) (the list
\(x_p,\ldots,x_{p+q-1}\));
\(\mathrm{At}(n,i,L)=N_{16}(N_{15}(n,i),L)\) (the atom \(X^n_i\) applied to
the term list \(L\));
\(\mathrm{All}(p,0,\varphi)=\varphi\),
\(\mathrm{All}(p,q+1,\varphi)=\mathrm{All}(p,q,N_{11}(p+q,\varphi))\), so that
\(\mathrm{All}(p,q,\ulcorner\varphi\urcorner)\) codes
\(\forall x_p\cdots\forall x_{p+q-1}\varphi\), and \(\mathrm{Ex}\) likewise with
\(N_{12}\);
\(\mathrm{Iff}(\varphi,\psi)=N_8(N_{10}(\varphi,\psi),N_{10}(\psi,\varphi))\);
and \(\mathrm{Cj}(L,\varphi)\), the right-grouped conjunction of the entries of
\(L\) followed by \(\varphi\), computed by a loop over the reversed list with
update \(\psi\mapsto N_8(e,\psi)\) starting from \(\varphi\).
Put \(X(c)=N_{15}(\mathrm{ar}(c)+1,c)\), the code of \(X_\sigma\) when
\(c=\ulcorner\sigma\urcorner\), and \(\mathrm{Lhs}(c,L)=N_{16}(X(c),L)\).

\emph{Step 4: the graph condition.}  Define \(\mathrm{Df}(c)\) by cases on the
tag of \(c\), with \(k=\mathrm{ar}(c)\) and \(\vec v=V(1,k)\):
\begin{itemize}
 \item tags \(30\), \(31\), \(32\): \(\mathrm{All}(1,k,N_{11}(0,\mathrm{Iff}(\mathrm{Lhs}(c,\vec v\mathbin{+\!\!+}[v(0)]),N_5(v(0),e))))\)
   with \(e=N_1()\), \(e=N_2(v(1))\), \(e=v((a)_1)\) respectively;
 \item tag \(33\), with \(m=\ell(L)\) and \(g=(a)_1\): the list
   \(B=[\mathrm{At}(k+1,(L)_0,\vec v\mathbin{+\!\!+}[v(k+1)]),\ldots,\mathrm{At}(k+1,(L)_{m-1},\vec v\mathbin{+\!\!+}[v(k+m)])]\),
   produced by a loop of length \(m\); then
   \(\mathrm{All}(1,k,N_{11}(0,\mathrm{Iff}(\mathrm{Lhs}(c,\vec v\mathbin{+\!\!+}[v(0)]),
   \mathrm{Ex}(k+1,m,\mathrm{Cj}(B,\mathrm{At}(m+1,g,V(k+1,m)\mathbin{+\!\!+}[v(0)]))))))\);
 \item tag \(34\), with \(g=(a)_0\), \(e=(a)_1\) and \(k'=\mathrm{ar}(g)\), so
   \(\vec w=V(1,k')\): the code \(N_8(D_0,D_1)\), where
   \(D_0=\mathrm{All}(1,k',N_{11}(0,\mathrm{Iff}(\mathrm{Lhs}(c,\vec w\mathbin{+\!\!+}[N_1(),v(0)]),\mathrm{At}(k'+1,g,\vec w\mathbin{+\!\!+}[v(0)]))))\)
   and
   \(D_1=\mathrm{All}(1,k'+1,N_{11}(0,\mathrm{Iff}(\mathrm{Lhs}(c,\vec w\mathbin{+\!\!+}[N_2(v(k'+1)),v(0)]),
   N_{12}(k'+2,N_8(\mathrm{Lhs}(c,\vec w\mathbin{+\!\!+}[v(k'+1),v(k'+2)]),\mathrm{At}(k'+3,e,\vec w\mathbin{+\!\!+}[v(k'+1),v(k'+2),v(0)]))))))\);
 \item tag \(35\), with \(g=(a)_0\): the code of the clause of
   \cref{def:p02-w02-p02-041-graph-sentence} item 6 assembled in the same way,
   whose inner atoms are
   \(\mathrm{At}(k+2,g,\vec v\mathbin{+\!\!+}[v(0),N_1()])\),
   \(N_5(N_3(v(k+1),N_2(v(k+2))),v(0))\),
   \(\mathrm{At}(k+2,g,\vec v\mathbin{+\!\!+}[v(k+1),v(k+3)])\) and
   \(N_7(N_5(v(k+3),N_1()))\), combined with \(N_8\), \(N_{10}\),
   \(N_{11}\), \(N_{12}\) and \(\mathrm{Iff}\) in the displayed order.
\end{itemize}
Each case is a fixed composition of primitive recursive functions with
\(\mathrm{ar}\) and list access, so \(\mathrm{Df}\) is primitive recursive.
Comparing each case with \cref{def:p02-w02-p02-041-graph-sentence} and with the
constructor table and \cref{def:p02-w02-p02-040-so-coding} shows that
\(\mathrm{Df}(\ulcorner\sigma\urcorner)=\ulcorner\mathsf{Def}_\sigma\urcorner\) for
every term \(\sigma\).  (In the tag \(34\) case \(\mathrm{ar}(c)=k'+1\), so
\(X(c)\) has the arity \(k'+2\) used by its atoms; in the tag \(33\) case
\(\mathrm{ar}(g)=m\).)

\emph{Step 5: assembly.}  Let \(p\) be the fixed code of the sentence
\(\mathsf{PA}_2\).  For a list \(L\) and a code \(\varphi\), let
\(\mathrm{Pf}(L,\varphi)\) be computed by a loop over the reversed list with
update \(\psi\mapsto N_{17}(\mathrm{ar}(e)+1,e,\psi)\) for each entry \(e\),
starting from \(\varphi\); it codes the prefix
\(\forall X_{\sigma_1}\cdots\forall X_{\sigma_n}\) in front of \(\varphi\) when
\(L=[\ulcorner\sigma_1\urcorner,\ldots,\ulcorner\sigma_n\urcorner]\).  Let
\(\mathrm{Mp}(L)\) be the list obtained by applying \(\mathrm{Df}\) to each
entry, by a loop of length \(\ell(L)\).  With \(S=\mathrm{Sb}(c)\), whose
length is at least \(1\) for a term code, let \(\mathrm{Cj}'(S)\) be
\(\mathrm{Cj}(L',\mathrm{Df}((S)_{\ell(S)-1}))\), where \(L'\) is
\(\mathrm{Mp}(S)\) with its last entry removed.  Define
\[
 \mathrm{rd}(c)=\begin{cases}
 N_{10}\Bigl(p,\ \mathrm{Pf}\bigl(S,\ N_{10}\bigl(\mathrm{Cj}'(S),\
   N_7(N_{12}(0,\mathrm{At}(2,c,[\nu(c),v(0)])))\bigr)\bigr)\Bigr),
 &\mathrm{RT}(c)=1\text{ and }\mathrm{ar}(c)=1,\\
 N_6(),&\text{otherwise}.
 \end{cases}
\]
This is a definition by cases from primitive recursive functions, so
\(\mathrm{rd}\) is primitive recursive.  For a term \(\rho\) of arity \(1\)
with code \(c\), Steps 1, 2 and 4 show that the first case applies and that
its value codes
\(\mathsf{PA}_2\to\forall X_{\sigma_1}\cdots\forall X_{\sigma_n}((\mathsf{Def}_{\sigma_1}\land\cdots\land\mathsf{Def}_{\sigma_n})\to\neg\exists x_0\,X_\rho(\overline c,x_0))\),
that is, \(\ulcorner\Theta_\rho\urcorner\); here \(X(c)=N_{15}(2,c)\) is the code of
\(X_\rho\).  On every other input the second case applies.
\end{proof}
\end{annexitem}

\begin{annexitem}{annex:p02-w02-p02-041-no-recursive-domain}{thm:p02-w02-p02-041-no-recursive-domain}{No recursive enumeration of full second-order validity}
\begin{proof}
\emph{(a)}  Suppose \(\beta\) is a term of arity \(1\) with
\(\operatorname{dom}(f_\beta)=\mathrm{Val}_2\).  By
\cref{thm:p02-w02-p02-041-effective} the function \(\mathrm{rd}\) is
primitive recursive, so \cref{thm:p02-w02-p02-041-coverage}(b) gives a term
\(\alpha\) of arity \(1\) with \(f_\alpha=\mathrm{rd}\).  Put
\(\delta=\mathbf C^1(\beta;\alpha)\); it is a term of arity \(1\), because
\(\mathrm{ar}(\beta)=1\) is the number of inner terms and
\(\mathrm{ar}(\alpha)=1\).  Since \(\mathrm{rd}\) is total,
\(f_\delta(x)\) is defined iff \(\mathrm{rd}(x)\in\operatorname{dom}(f_\beta)=\mathrm{Val}_2\).
Let \(d=\ulcorner\delta\urcorner\).  As \(\delta\) is a term of arity \(1\),
\(\mathrm{rd}(d)=\ulcorner\Theta_\delta\urcorner\), the code of a second-order
sentence, so \(\mathrm{rd}(d)\in\mathrm{Val}_2\) iff
\(\entails_{\mathrm{full}}\Theta_\delta\).  By
\cref{thm:p02-w02-p02-041-reduction} this holds iff \(f_\delta(d)\) is
undefined.  Together: \(f_\delta(d)\) is defined iff it is undefined, a
contradiction.  The second sentence of (a) follows by
\cref{thm:p02-w02-p02-041-coverage}(a).

\emph{(b)}  Suppose \(\mathsf D\) is such a calculus, with theorem codes
\(\mathrm{Th}=\operatorname{dom}(\chi)\) for a partial recursive \(\chi\).  Let
\(\mathrm{Sent}_2\) be the characteristic function of the set of codes of
second-order sentences of the signature; it is primitive recursive by the
second-order recognition clauses in the proof of
\cref{thm:p02-w02-p02-040-no-effective-calculus}
(\cref{annex:p02-w02-p02-040-no-effective-calculus}).  The hypothesis on
\(\mathsf D\) says that a sentence code is in \(\mathrm{Th}\) iff it is in
\(\mathrm{Val}_2\), and every member of \(\mathrm{Val}_2\) is a sentence code;
so \(\mathrm{Val}_2=\mathrm{Th}\cap\{c:\mathrm{Sent}_2(c)=1\}\).  The map
\(g(c,y)=1\mathbin{dot{-}}\mathrm{Sent}_2(c)\) is primitive recursive (a composition
of truncated subtraction, the constant \(1=S(Z_2(c,y))\), \(\mathrm{Sent}_2\)
and a projection), so by \cref{thm:p02-w02-p02-041-coverage} there are terms
\(\kappa\) of arity \(1\) with \(f_\kappa=\chi\) and \(\zeta\) of arity \(2\)
with \(f_\zeta=g\).  The term \(\mathbf M(\zeta)\) has arity \(1\):
if \(\mathrm{Sent}_2(c)=1\) then \(f_\zeta(c,0)=0\) and
\(f_{\mathbf M(\zeta)}(c)=0\); if \(\mathrm{Sent}_2(c)=0\) then
\(f_\zeta(c,y)=1\) for every \(y\), no zero is reached, and
\(f_{\mathbf M(\zeta)}(c)\) is undefined.  Let
\(\beta=\mathbf C^1(\mathbf P^2_1;\kappa,\mathbf M(\zeta))\), a term of arity
\(1\).  Then \(f_\beta(c)\) is defined iff \(\chi(c)\) and
\(f_{\mathbf M(\zeta)}(c)\) are both defined, iff
\(c\in\mathrm{Th}\) and \(\mathrm{Sent}_2(c)=1\), iff \(c\in\mathrm{Val}_2\).
This contradicts (a).

\emph{(c)}  If \(\mathrm{Val}_2\) were computably enumerable, it would be the
domain of a partial computable \(\psi\) (\cref{sec:turing-machines-03}); by
\cref{thm:p04-tm-partial-recursive} \(\psi\) is partial recursive, contrary to
(a).

\emph{(d)}  If the theorem codes of such a calculus formed a computably
enumerable set, they would be the domain of a partial computable function,
which is partial recursive by \cref{thm:p04-tm-partial-recursive}, contrary to
(b).
\end{proof}
\end{annexitem}

\section{Proofs for truth preservation through encoding and decoding}
\label{sec:annex-p02-w02-p02-041-decoding}

Notation and conventions are those of
\cref{sec:p02-w02-p02-041-decoding}.  For a matching assignment \(s\), an
object variable \(x\), an element \(a\) of the object domain, a relation
variable \(X^n\) and a relation \(R\) in the relation domain, \(s[x\mapsto a]\)
and \(s[X^n\mapsto R]\) are again matching assignments; the same holds for
sorted assignments of a many-sorted structure and elements of the
corresponding sort.

\begin{annexitem}{annex:p02-w02-p02-041-henkin-coincidence}{thm:p02-w02-p02-041-henkin-coincidence}{Coincidence for Henkin semantics}
\begin{proof}
\emph{Terms.}  If \(s\) and \(t\) agree on the variables of an object term
\(r\), then \(\operatorname{val}^{\mathcal H}_s(r)=\operatorname{val}^{\mathcal
H}_t(r)\): both are first-order values in the object reduct of
\(\mathcal H\) under the restrictions of \(s\) and \(t\) to object variables,
and these restrictions agree on the variables of \(r\), so
the term-coincidence lemma of \cref{sec:fol-semantics-01} applies.

\emph{Formulas.}  We prove, by induction on \(\varphi\), that for all matching
assignments \(s,t\) agreeing on the free variables of \(\varphi\),
\(\mathcal H,s\entails\varphi\) iff \(\mathcal H,t\entails\varphi\).

\emph{Atoms.}  For \(r=u\), the variables of \(r\) and \(u\) are free in the
atom, so the two values of \(r\) agree and the two values of \(u\) agree; both
sides ask whether the same two elements are equal.  For \(P(t_1,\ldots,t_n)\)
the values of the \(t_i\) agree, and both sides ask whether the same tuple
belongs to \(P^{\mathcal H}\).  For \(X^n(t_1,\ldots,t_n)\) the variable
\(X^n\) is free in the atom, so \(s(X^n)=t(X^n)\); the values of the \(t_i\)
agree, and both sides ask whether the same tuple belongs to the same relation.
The atom \(\bottom\) is false under both assignments.

\emph{Connectives.}  For \(\neg\psi\), the free variables of \(\psi\) are
those of \(\neg\psi\); by the induction hypothesis \(\psi\) has the same truth
value under \(s\) and \(t\), hence so does \(\neg\psi\).  For
\(\psi\circ\chi\) with \(\circ\) one of \(\land,\lor,\to,\leftrightarrow\),
the free variables of \(\psi\) and those of \(\chi\) are free in
\(\psi\circ\chi\), so \(s,t\) agree on each set; by the induction hypothesis
\(\psi\) has one truth value under both assignments, and so has \(\chi\).  The
clause for \(\circ\) computes the truth value of \(\psi\circ\chi\) from the
pair of these two values alone, so \(\psi\circ\chi\) has one truth value under
both assignments.

\emph{Object quantifiers.}  Let \(\varphi\) be \(\forall x\,\psi\) or
\(\exists x\,\psi\).  Every free variable of \(\psi\) is \(x\) or is free in
\(\varphi\).  For \(a\in M\), the assignments \(s[x\mapsto a]\) and
\(t[x\mapsto a]\) both send \(x\) to \(a\) and agree with \(s\), \(t\)
respectively, hence with each other, on the remaining free variables of
\(\psi\).  By the induction hypothesis the sets
\(A_s=\{a\in M:\mathcal H,s[x\mapsto a]\entails\psi\}\) and
\(A_t=\{a\in M:\mathcal H,t[x\mapsto a]\entails\psi\}\) are equal.  The
universal clause holds under \(s\) iff \(A_s=M\), and the existential clause
iff \(A_s\ne\varnothing\); likewise for \(t\) with \(A_t\).  Since
\(A_s=A_t\), \(\varphi\) has the same truth value under \(s\) and \(t\).

\emph{Relation quantifiers.}  Let \(\varphi\) be \(\forall X^n\,\psi\) or
\(\exists X^n\,\psi\).  For \(R\in\mathcal R_n\), the assignments
\(s[X^n\mapsto R]\) and \(t[X^n\mapsto R]\) are matching, send \(X^n\) to
\(R\), and agree on the other free variables of \(\psi\), all free in
\(\varphi\).  By the induction hypothesis
\(\{R\in\mathcal R_n:\mathcal H,s[X^n\mapsto R]\entails\psi\}=
\{R\in\mathcal R_n:\mathcal H,t[X^n\mapsto R]\entails\psi\}\).  The Henkin
clauses for \(\forall X^n\) and \(\exists X^n\) ask whether this set is all of
\(\mathcal R_n\) or is nonempty, so \(\varphi\) has the same truth value under
\(s\) and \(t\).

A sentence has no free variables, so any two matching assignments agree on
them, which gives the final statement.
\end{proof}
\end{annexitem}

\begin{annexitem}{annex:p02-w02-p02-041-tagging}{thm:p02-w02-p02-041-tagging}{Tagging invariant for the incidence encoding}
\begin{proof}
\emph{(a)}  \(D^{\mathcal H^\sharp}=\{D\}\times M\) is nonempty because
\(M\) is, and \(S_n^{\mathcal H^\sharp}=\{S_n\}\times\mathcal R_n\) is
nonempty because \(\mathcal R_n\) is.  Members of different sorts have
different first components, so the domains are pairwise disjoint.  Each
\(c^{\mathcal H^\sharp}=(D,c^{\mathcal H})\) lies in
\(D^{\mathcal H^\sharp}\), and each \(f^{\mathcal H^\sharp}\) maps
\((D^{\mathcal H^\sharp})^k\) into \(D^{\mathcal H^\sharp}\).  By definition
\(E_n^{\mathcal H^\sharp}\subseteq
S_n^{\mathcal H^\sharp}\times(D^{\mathcal H^\sharp})^n\).  The constant
\(\widehat P:S_n\) receives \((S_n,P^{\mathcal H})\), a member of
\(S_n^{\mathcal H^\sharp}\) because \(P^{\mathcal H}\in\mathcal R_n\); this
is the only use of that hypothesis.  Finally
\(s^\sharp(x)=(D,s(x))\in D^{\mathcal H^\sharp}\) and
\(s^\sharp(X^n)=(S_n,s(X^n))\in S_n^{\mathcal H^\sharp}\), since
\(s(X^n)\in\mathcal R_n\).

\emph{(b)}  Induction on \(t\).  For a variable,
\(\operatorname{val}^{\mathcal H^\sharp}_{s^\sharp}(x)=s^\sharp(x)=(D,s(x))\).
For a constant, the value is \(c^{\mathcal H^\sharp}=(D,c^{\mathcal H})\).
For \(f(t_1,\ldots,t_k)\), with \(a_i=\operatorname{val}^{\mathcal H}_s(t_i)\),
the induction hypothesis gives
\[
\begin{aligned}
 \operatorname{val}^{\mathcal H^\sharp}_{s^\sharp}(f(t_1,\ldots,t_k))
 &=f^{\mathcal H^\sharp}((D,a_1),\ldots,(D,a_k))\\
 &=(D,f^{\mathcal H}(a_1,\ldots,a_k))
 =(D,\operatorname{val}^{\mathcal H}_s(f(t_1,\ldots,t_k))).
\end{aligned}
\]

\emph{Update identities.}  For \(a\in M\) and \(R\in\mathcal R_n\),
\[
 (s[x\mapsto a])^\sharp=s^\sharp[x\mapsto(D,a)],
 \qquad
 (s[X^n\mapsto R])^\sharp=s^\sharp[X^n\mapsto(S_n,R)],
\]
because both sides take the displayed value at the updated variable, by the
definition of \(\sharp\) on assignments, and agree with \(s^\sharp\) at every
other variable.

\emph{(c)}  Induction on \(\varphi\); at each stage the claim is proved for
every matching assignment \(s\).  In the atomic cases write
\(a_i=\operatorname{val}^{\mathcal H}_s(t_i)\); by (b) the value of \(t_i\) in
\(\mathcal H^\sharp\) under \(s^\sharp\) is \((D,a_i)\).

\emph{Equality.}  \((r=u)^\sharp\) is \(r=u\) at sort \(D\).  By (b),
\(\mathcal H^\sharp,s^\sharp\entails r=u\) iff
\((D,\operatorname{val}^{\mathcal H}_s(r))=(D,\operatorname{val}^{\mathcal
H}_s(u))\) iff \(\operatorname{val}^{\mathcal H}_s(r)=\operatorname{val}^{\mathcal
H}_s(u)\) iff \(\mathcal H,s\entails r=u\).

\emph{Relation-variable atom.}  \((X^n(t_1,\ldots,t_n))^\sharp\) is
\(E_n(X^n,t_1,\ldots,t_n)\), and
\[
\begin{aligned}
 \mathcal H^\sharp,s^\sharp\entails E_n(X^n,\bar t)
 &\iff\bigl((S_n,s(X^n)),(D,a_1),\ldots,(D,a_n)\bigr)\in E_n^{\mathcal H^\sharp}\\
 &\iff(a_1,\ldots,a_n)\in s(X^n)
 \iff\mathcal H,s\entails X^n(\bar t),
\end{aligned}
\]
by the definition of \(E_n^{\mathcal H^\sharp}\) and the Henkin atomic clause
(the clause of \cref{sec:second-order-logic-02} with \(s(X^n)\in\mathcal
R_n\)).

\emph{Predicate-constant atom.}  \((P(t_1,\ldots,t_n))^\sharp\) is
\(E_n(\widehat P,t_1,\ldots,t_n)\), and
\[
\begin{aligned}
 \mathcal H^\sharp,s^\sharp\entails E_n(\widehat P,\bar t)
 &\iff\bigl((S_n,P^{\mathcal H}),(D,a_1),\ldots,(D,a_n)\bigr)\in E_n^{\mathcal H^\sharp}\\
 &\iff(a_1,\ldots,a_n)\in P^{\mathcal H}
 \iff\mathcal H,s\entails P(\bar t).
\end{aligned}
\]

\emph{Falsum.}  \(\bottom^\sharp=\bottom\), false on both sides.

\emph{Negation.}  \(\mathcal H^\sharp,s^\sharp\entails(\neg\psi)^\sharp\) iff
\(\mathcal H^\sharp,s^\sharp\entails\neg\psi^\sharp\) iff
\(\mathcal H^\sharp,s^\sharp\not\entails\psi^\sharp\) iff, by the induction
hypothesis, \(\mathcal H,s\not\entails\psi\) iff
\(\mathcal H,s\entails\neg\psi\).

\emph{Conjunction.}  \(\mathcal H^\sharp,s^\sharp\entails(\psi\land\chi)^\sharp\)
iff \(\mathcal H^\sharp,s^\sharp\entails\psi^\sharp\) and
\(\mathcal H^\sharp,s^\sharp\entails\chi^\sharp\) iff, by the induction
hypothesis for \(\psi\) and for \(\chi\), \(\mathcal H,s\entails\psi\) and
\(\mathcal H,s\entails\chi\) iff \(\mathcal H,s\entails\psi\land\chi\).

\emph{Disjunction.}  \(\mathcal H^\sharp,s^\sharp\entails(\psi\lor\chi)^\sharp\)
iff \(\mathcal H^\sharp,s^\sharp\entails\psi^\sharp\) or
\(\mathcal H^\sharp,s^\sharp\entails\chi^\sharp\) iff, by the induction
hypothesis, \(\mathcal H,s\entails\psi\) or \(\mathcal H,s\entails\chi\) iff
\(\mathcal H,s\entails\psi\lor\chi\).

\emph{Implication.}  \(\mathcal H^\sharp,s^\sharp\entails(\psi\to\chi)^\sharp\)
iff \(\mathcal H^\sharp,s^\sharp\not\entails\psi^\sharp\) or
\(\mathcal H^\sharp,s^\sharp\entails\chi^\sharp\) iff, by the induction
hypothesis, \(\mathcal H,s\not\entails\psi\) or \(\mathcal H,s\entails\chi\)
iff \(\mathcal H,s\entails\psi\to\chi\).

\emph{Biconditional.}
\(\mathcal H^\sharp,s^\sharp\entails(\psi\leftrightarrow\chi)^\sharp\) iff
\(\psi^\sharp\) and \(\chi^\sharp\) have the same truth value in
\(\mathcal H^\sharp\) under \(s^\sharp\) iff, by the induction hypothesis,
\(\psi\) and \(\chi\) have the same truth value in \(\mathcal H\) under \(s\)
iff \(\mathcal H,s\entails\psi\leftrightarrow\chi\).

\emph{Universal object quantifier.}  The map \(a\mapsto(D,a)\) is a bijection
from \(M\) onto \(D^{\mathcal H^\sharp}\).  Hence
\[
\begin{aligned}
 \mathcal H^\sharp,s^\sharp\entails\forall x{:}D\,\psi^\sharp
 &\iff\text{for every }b\in D^{\mathcal H^\sharp},\
   \mathcal H^\sharp,s^\sharp[x\mapsto b]\entails\psi^\sharp\\
 &\iff\text{for every }a\in M,\
   \mathcal H^\sharp,s^\sharp[x\mapsto(D,a)]\entails\psi^\sharp\\
 &\iff\text{for every }a\in M,\
   \mathcal H^\sharp,(s[x\mapsto a])^\sharp\entails\psi^\sharp\\
 &\iff\text{for every }a\in M,\
   \mathcal H,s[x\mapsto a]\entails\psi\\
 &\iff\mathcal H,s\entails\forall x\,\psi,
\end{aligned}
\]
using the bijection, the update identity, and the induction hypothesis at the
matching assignment \(s[x\mapsto a]\).

\emph{Existential object quantifier.}  The same four steps with ``for some''
in place of ``for every'':
\(\mathcal H^\sharp,s^\sharp\entails\exists x{:}D\,\psi^\sharp\) iff some
\(b\in D^{\mathcal H^\sharp}\) has
\(\mathcal H^\sharp,s^\sharp[x\mapsto b]\entails\psi^\sharp\) iff some
\(a\in M\) has \(\mathcal H^\sharp,(s[x\mapsto a])^\sharp\entails\psi^\sharp\)
(each \(b\) is \((D,a)\) for one \(a\)) iff some \(a\in M\) has
\(\mathcal H,s[x\mapsto a]\entails\psi\) iff \(\mathcal H,s\entails\exists
x\,\psi\).

\emph{Universal relation quantifier.}  The map \(R\mapsto(S_n,R)\) is a
bijection from \(\mathcal R_n\) onto \(S_n^{\mathcal H^\sharp}\).  Hence
\[
\begin{aligned}
 \mathcal H^\sharp,s^\sharp\entails\forall X^n{:}S_n\,\psi^\sharp
 &\iff\text{for every }R\in\mathcal R_n,\
   \mathcal H^\sharp,s^\sharp[X^n\mapsto(S_n,R)]\entails\psi^\sharp\\
 &\iff\text{for every }R\in\mathcal R_n,\
   \mathcal H^\sharp,(s[X^n\mapsto R])^\sharp\entails\psi^\sharp\\
 &\iff\text{for every }R\in\mathcal R_n,\
   \mathcal H,s[X^n\mapsto R]\entails\psi\\
 &\iff\mathcal H,s\entails\forall X^n\,\psi,
\end{aligned}
\]
the last step being the Henkin clause, under which \(X^n\) ranges over
\(\mathcal R_n\).

\emph{Existential relation quantifier.}
\(\mathcal H^\sharp,s^\sharp\entails\exists X^n{:}S_n\,\psi^\sharp\) iff some
member of \(S_n^{\mathcal H^\sharp}\), necessarily \((S_n,R)\) for one
\(R\in\mathcal R_n\), satisfies the matrix, iff some \(R\in\mathcal R_n\) has
\(\mathcal H^\sharp,(s[X^n\mapsto R])^\sharp\entails\psi^\sharp\) iff some
\(R\in\mathcal R_n\) has \(\mathcal H,s[X^n\mapsto R]\entails\psi\) iff
\(\mathcal H,s\entails\exists X^n\,\psi\).

These cases exhaust the formula clauses of
\cref{sec:second-order-logic-01} and the translation clauses of
\cref{def:p02-sol-incidence-translation}.

\emph{(d)}  Let \(w\) be a sorted assignment of \(\mathcal H^\sharp\).  For
each object variable \(x\), \(w(x)\in\{D\}\times M\), so \(w(x)=(D,a_x)\) for
exactly one \(a_x\in M\); for each relation variable \(X^n\),
\(w(X^n)=(S_n,R_X)\) for exactly one \(R_X\in\mathcal R_n\).  Put
\(u(x)=a_x\) and \(u(X^n)=R_X\).  Then \(u\) is a matching assignment and
\(u^\sharp=w\).  If also \(u'^\sharp=w\), then \((D,u'(x))=(D,u(x))\) and
\((S_n,u'(X^n))=(S_n,u(X^n))\) for every variable, so \(u'=u\).  Each value
of \(u\) is the second component of a given pair, so no choice is made.

\emph{(e)}  By \cref{def:p02-sol-derived-calculus},
\(T_{\mathsf{CC}}^\sharp=\{\theta^\sharp:\theta\text{ an eligible }
\mathsf{Comp}^n_\chi\text{ or }\mathsf{Choice}^n_\chi\text{ instance}\}\).
Since the predicate-constant condition holds by hypothesis, \(\mathcal H\) is
an \(\mathsf H_{\mathsf{CC}}\)-model iff \(\mathcal H,u\entails\theta\) for
every eligible \(\theta\) and every matching \(u\); by (c) iff
\(\mathcal H^\sharp,u^\sharp\entails\theta^\sharp\) for every eligible
\(\theta\) and every matching \(u\); and by (d), since \(u\mapsto u^\sharp\)
is onto the sorted assignments of \(\mathcal H^\sharp\), iff
\(\mathcal H^\sharp,w\entails\theta^\sharp\) for every eligible \(\theta\) and
every sorted assignment \(w\).
\end{proof}
\end{annexitem}

\begin{annexitem}{annex:p02-w02-p02-041-collapse}{thm:p02-w02-p02-041-collapse}{Truth preservation under extensional collapse}
\begin{proof}
\emph{(a)}  \(M=D^{\mathcal A}\) is nonempty because the domains of a strict
structure are nonempty.  Each \(\operatorname{ext}_n(r)\) is a subset of
\((D^{\mathcal A})^n=M^n\), so \(\mathcal R_n\subseteq\mathcal P(M^n)\), and
\(\mathcal R_n\) is nonempty because \(S_n^{\mathcal A}\) is.  Object
constants and function symbols are declared at sort \(D\) in the incidence
signature, so \(c^{\mathcal A}\in M\) and \(f^{\mathcal A}\) maps \(M^k\) into
\(M\); the object reduct is a first-order structure with domain \(M\).  Since
\(\widehat P:S_n\), \(\widehat P^{\mathcal A}\in S_n^{\mathcal A}\) and
\(P^{\mathcal H_{\mathrm{ext}}}=\operatorname{ext}_n(\widehat P^{\mathcal
A})\in\mathcal R_n\).  Thus \(\mathcal H_{\mathrm{ext}}\) is a Henkin
structure as defined in \cref{sec:second-order-logic-05}.  Finally
\(s^{\mathrm{ext}}(x)=s(x)\in D^{\mathcal A}=M\), and
\(s^{\mathrm{ext}}(X^n)=\operatorname{ext}_n(s(X^n))\in\mathcal R_n\) because
\(s(X^n)\in S_n^{\mathcal A}\); so \(s^{\mathrm{ext}}\) is matching.

\emph{(b)}  Induction on \(t\).  For a variable,
\(s^{\mathrm{ext}}(x)=s(x)\).  For a constant,
\(c^{\mathcal H_{\mathrm{ext}}}=c^{\mathcal A}\).  For
\(f(t_1,\ldots,t_k)\), the function \(f^{\mathcal H_{\mathrm{ext}}}=f^{\mathcal
A}\) is applied to arguments that are equal by the induction hypothesis.

\emph{Update identities.}  For \(a\in D^{\mathcal A}\) and
\(r\in S_n^{\mathcal A}\),
\[
 (s[x\mapsto a])^{\mathrm{ext}}=s^{\mathrm{ext}}[x\mapsto a],
 \qquad
 (s[X^n\mapsto r])^{\mathrm{ext}}=s^{\mathrm{ext}}[X^n\mapsto\operatorname{ext}_n(r)],
\]
because both sides take the displayed value at the updated variable, by the
definition of the decoded assignment, and agree with \(s^{\mathrm{ext}}\)
elsewhere.

\emph{(c)}  Induction on \(\varphi\); at each stage the claim is proved for
every sorted assignment \(s\) of \(\mathcal A\).  In the atomic cases write
\(a_i=\operatorname{val}^{\mathcal A}_s(t_i)\), which by (b) is also the
value of \(t_i\) in \(\mathcal H_{\mathrm{ext}}\) under \(s^{\mathrm{ext}}\);
each \(a_i\) lies in \(D^{\mathcal A}\) because \(t_i\) has sort \(D\).

\emph{Equality.}  \((r=u)^\sharp\) is \(r=u\).
\(\mathcal A,s\entails r=u\) iff
\(\operatorname{val}^{\mathcal A}_s(r)=\operatorname{val}^{\mathcal A}_s(u)\)
iff, by (b), the two values in \(\mathcal H_{\mathrm{ext}}\) under
\(s^{\mathrm{ext}}\) are equal, iff
\(\mathcal H_{\mathrm{ext}},s^{\mathrm{ext}}\entails r=u\).

\emph{Incidence atom.}  \((X^n(t_1,\ldots,t_n))^\sharp\) is
\(E_n(X^n,t_1,\ldots,t_n)\), and
\[
\begin{aligned}
 \mathcal A,s\entails E_n(X^n,\bar t)
 &\iff(s(X^n),a_1,\ldots,a_n)\in E_n^{\mathcal A}\\
 &\iff(a_1,\ldots,a_n)\in\operatorname{ext}_n(s(X^n))\\
 &\iff(a_1,\ldots,a_n)\in s^{\mathrm{ext}}(X^n)\\
 &\iff\mathcal H_{\mathrm{ext}},s^{\mathrm{ext}}\entails X^n(\bar t).
\end{aligned}
\]
The second step is the definition of \(\operatorname{ext}_n\), applicable
because \((a_1,\ldots,a_n)\in(D^{\mathcal A})^n\); the third is the
definition of \(s^{\mathrm{ext}}\); the fourth is the Henkin atomic clause.

\emph{Predicate-constant atom.}  \((P(t_1,\ldots,t_n))^\sharp\) is
\(E_n(\widehat P,t_1,\ldots,t_n)\), and
\[
\begin{aligned}
 \mathcal A,s\entails E_n(\widehat P,\bar t)
 &\iff(\widehat P^{\mathcal A},a_1,\ldots,a_n)\in E_n^{\mathcal A}\\
 &\iff(a_1,\ldots,a_n)\in\operatorname{ext}_n(\widehat P^{\mathcal A})
   =P^{\mathcal H_{\mathrm{ext}}}\\
 &\iff\mathcal H_{\mathrm{ext}},s^{\mathrm{ext}}\entails P(\bar t).
\end{aligned}
\]

\emph{Falsum.}  \(\bottom^\sharp=\bottom\), false on both sides.

\emph{Negation.}  \(\mathcal A,s\entails\neg\psi^\sharp\) iff
\(\mathcal A,s\not\entails\psi^\sharp\) iff, by the induction hypothesis,
\(\mathcal H_{\mathrm{ext}},s^{\mathrm{ext}}\not\entails\psi\) iff
\(\mathcal H_{\mathrm{ext}},s^{\mathrm{ext}}\entails\neg\psi\).

\emph{Conjunction.}  \(\mathcal A,s\entails\psi^\sharp\land\chi^\sharp\) iff
\(\mathcal A,s\entails\psi^\sharp\) and \(\mathcal A,s\entails\chi^\sharp\)
iff, by the induction hypothesis for \(\psi\) and for \(\chi\),
\(\mathcal H_{\mathrm{ext}},s^{\mathrm{ext}}\entails\psi\) and
\(\mathcal H_{\mathrm{ext}},s^{\mathrm{ext}}\entails\chi\) iff
\(\mathcal H_{\mathrm{ext}},s^{\mathrm{ext}}\entails\psi\land\chi\).

\emph{Disjunction.}  \(\mathcal A,s\entails\psi^\sharp\lor\chi^\sharp\) iff
\(\mathcal A,s\entails\psi^\sharp\) or \(\mathcal A,s\entails\chi^\sharp\)
iff, by the induction hypothesis,
\(\mathcal H_{\mathrm{ext}},s^{\mathrm{ext}}\entails\psi\) or
\(\mathcal H_{\mathrm{ext}},s^{\mathrm{ext}}\entails\chi\) iff
\(\mathcal H_{\mathrm{ext}},s^{\mathrm{ext}}\entails\psi\lor\chi\).

\emph{Implication.}  \(\mathcal A,s\entails\psi^\sharp\to\chi^\sharp\) iff
\(\mathcal A,s\not\entails\psi^\sharp\) or \(\mathcal A,s\entails\chi^\sharp\)
iff, by the induction hypothesis,
\(\mathcal H_{\mathrm{ext}},s^{\mathrm{ext}}\not\entails\psi\) or
\(\mathcal H_{\mathrm{ext}},s^{\mathrm{ext}}\entails\chi\) iff
\(\mathcal H_{\mathrm{ext}},s^{\mathrm{ext}}\entails\psi\to\chi\).

\emph{Biconditional.}
\(\mathcal A,s\entails\psi^\sharp\leftrightarrow\chi^\sharp\) iff
\(\psi^\sharp,\chi^\sharp\) have the same truth value in \(\mathcal A\) under
\(s\) iff, by the induction hypothesis, \(\psi,\chi\) have the same truth
value in \(\mathcal H_{\mathrm{ext}}\) under \(s^{\mathrm{ext}}\) iff
\(\mathcal H_{\mathrm{ext}},s^{\mathrm{ext}}\entails\psi\leftrightarrow\chi\).

\emph{Universal object quantifier.}  Since \(M=D^{\mathcal A}\),
\[
\begin{aligned}
 \mathcal A,s\entails\forall x{:}D\,\psi^\sharp
 &\iff\text{for every }a\in D^{\mathcal A},\
   \mathcal A,s[x\mapsto a]\entails\psi^\sharp\\
 &\iff\text{for every }a\in D^{\mathcal A},\
   \mathcal H_{\mathrm{ext}},(s[x\mapsto a])^{\mathrm{ext}}\entails\psi\\
 &\iff\text{for every }a\in M,\
   \mathcal H_{\mathrm{ext}},s^{\mathrm{ext}}[x\mapsto a]\entails\psi\\
 &\iff\mathcal H_{\mathrm{ext}},s^{\mathrm{ext}}\entails\forall x\,\psi,
\end{aligned}
\]
by the induction hypothesis at the sorted assignment \(s[x\mapsto a]\) and the
update identity.

\emph{Existential object quantifier.}
\(\mathcal A,s\entails\exists x{:}D\,\psi^\sharp\) iff some
\(a\in D^{\mathcal A}\) has \(\mathcal A,s[x\mapsto a]\entails\psi^\sharp\) iff,
by the induction hypothesis and the update identity, some \(a\in M\) has
\(\mathcal H_{\mathrm{ext}},s^{\mathrm{ext}}[x\mapsto a]\entails\psi\) iff
\(\mathcal H_{\mathrm{ext}},s^{\mathrm{ext}}\entails\exists x\,\psi\).

\emph{Universal relation quantifier.}  By the many-sorted clause, the
induction hypothesis at \(s[X^n\mapsto r]\), and the update identity,
\[
\begin{aligned}
 &\mathcal A,s\entails\forall X^n{:}S_n\,\psi^\sharp
 \quad\text{iff}\\
 &(\ast)\quad\text{for every }r\in S_n^{\mathcal A},\
 \mathcal H_{\mathrm{ext}},s^{\mathrm{ext}}[X^n\mapsto\operatorname{ext}_n(r)]\entails\psi .
\end{aligned}
\]
The Henkin clause says that
\(\mathcal H_{\mathrm{ext}},s^{\mathrm{ext}}\entails\forall X^n\,\psi\) iff
\((\ast\ast)\) for every \(R\in\mathcal R_n\),
\(\mathcal H_{\mathrm{ext}},s^{\mathrm{ext}}[X^n\mapsto R]\entails\psi\).
\((\ast\ast)\Rightarrow(\ast)\): for \(r\in S_n^{\mathcal A}\),
\(\operatorname{ext}_n(r)\in\mathcal R_n\), so \((\ast\ast)\) applies with
\(R=\operatorname{ext}_n(r)\).  \((\ast)\Rightarrow(\ast\ast)\): for
\(R\in\mathcal R_n\), the definition
\(\mathcal R_n=\operatorname{ext}_n[S_n^{\mathcal A}]\) provides some
\(r\in S_n^{\mathcal A}\) with \(\operatorname{ext}_n(r)=R\); \((\ast)\) at
that \(r\) gives \(\mathcal H_{\mathrm{ext}},s^{\mathrm{ext}}[X^n\mapsto
R]\entails\psi\).

\emph{Existential relation quantifier.}  Likewise by the many-sorted clause,
the induction hypothesis and the update identity,
\[
\begin{aligned}
 &\mathcal A,s\entails\exists X^n{:}S_n\,\psi^\sharp
 \quad\text{iff}\\
 &(\dagger)\quad\text{for some }r\in S_n^{\mathcal A},\
 \mathcal H_{\mathrm{ext}},s^{\mathrm{ext}}[X^n\mapsto\operatorname{ext}_n(r)]\entails\psi,
\end{aligned}
\]
and the Henkin clause says that
\(\mathcal H_{\mathrm{ext}},s^{\mathrm{ext}}\entails\exists X^n\,\psi\) iff
\((\dagger\dagger)\) for some \(R\in\mathcal R_n\),
\(\mathcal H_{\mathrm{ext}},s^{\mathrm{ext}}[X^n\mapsto R]\entails\psi\).
\((\dagger)\Rightarrow(\dagger\dagger)\): take
\(R=\operatorname{ext}_n(r)\in\mathcal R_n\).
\((\dagger\dagger)\Rightarrow(\dagger)\): the given \(R\) equals
\(\operatorname{ext}_n(r)\) for some \(r\in S_n^{\mathcal A}\), and that \(r\)
witnesses \((\dagger)\).

In both relation-quantifier cases only the set
\(\operatorname{ext}_n[S_n^{\mathcal A}]\) enters, so distinct objects with
equal extensions contribute the same instance, and each implication uses one
existential instantiation, so no choice principle is used.  The cases exhaust
the formula clauses and the translation clauses.

\emph{(d)}  The predicate-constant condition holds by (a).  Let \(\theta\) be
an eligible \(\mathsf{Comp}^n_\chi\) or \(\mathsf{Choice}^n_\chi\) instance.
In the notation of \cref{def:p02-sol-matched-henkin} the free variables of
\(\chi\) are among the displayed \(\bar x\), \(y\) (for choice), \(\bar p\)
and \(\bar Q\); the displayed quantifiers bind \(\bar x\), \(y\), \(z\),
\(X^n\) and \(F^{n+1}\), and \(\operatorname{ucl}\) binds \(\bar p,\bar Q\).
So \(\theta\) is a sentence.  As \(\theta^\sharp\in T_{\mathsf{CC}}^\sharp\),
the hypothesis gives \(\mathcal A,s\entails\theta^\sharp\); by (c),
\(\mathcal H_{\mathrm{ext}},s^{\mathrm{ext}}\entails\theta\); since
\(s^{\mathrm{ext}}\) is matching by (a),
\cref{thm:p02-w02-p02-041-henkin-coincidence} gives
\(\mathcal H_{\mathrm{ext}},t\entails\theta\) for every matching \(t\).  So
\(\mathcal H_{\mathrm{ext}}\) satisfies every eligible instance and is an
\(\mathsf H_{\mathsf{CC}}\)-model.
\end{proof}
\end{annexitem}

\begin{annexitem}{annex:p02-w02-p02-041-adequacy-route}{thm:p02-w02-p02-041-adequacy-route}{Matched adequacy through the encoding and the decoding}
\begin{proof}
Write \(\mathsf F\) for the flat one-sorted premise set
\(\mathsf{Ax}_{\mathrm{sort}}\cup\mathsf{VarGuard}_{\mathrm{MS}}\cup
\operatorname{tr}_{\mathrm{MS}}[T_{\mathsf{CC}}^\sharp\cup\Gamma^\sharp]\) of
\cref{def:p02-sol-derived-calculus}.  The incidence signature has countably
many sorts \(D,S_1,S_2,\ldots\), countably many symbols, and the countably
infinite variable stocks of the source language, so
\cref{thm:p02-many-sorted-translation} applies to it, and its flat signature is
countable.

\emph{Forward.}  Assume
\(T_{\mathsf{CC}}^\sharp\cup\Gamma^\sharp\derives_{\mathsf{MSFOL}}\varphi^\sharp\),
that is,
\(\mathsf F\derives_{\mathsf{ND}_{\mathsf{FOL}}}\operatorname{tr}_{\mathrm{MS}}(\varphi^\sharp)\).
By \cref{thm:fol-soundness},
\(\mathsf F\entails\operatorname{tr}_{\mathrm{MS}}(\varphi^\sharp)\), and by
\cref{thm:p02-many-sorted-translation},
\(T_{\mathsf{CC}}^\sharp\cup\Gamma^\sharp\entails_{\mathrm{MS}}\varphi^\sharp\).
Let \(\mathcal H\) be an \(\mathsf H_{\mathsf{CC}}\)-model and \(s\) a
matching assignment with \(\mathcal H,s\entails\psi\) for every
\(\psi\in\Gamma\).  By \cref{thm:p02-w02-p02-041-tagging}(a),
\(\mathcal H^\sharp\) is a strict structure and \(s^\sharp\) a sorted
assignment; by (c), \(\mathcal H^\sharp,s^\sharp\entails\psi^\sharp\) for
every \(\psi\in\Gamma\); by (e), \(\mathcal H^\sharp,s^\sharp\entails\theta\)
for every \(\theta\in T_{\mathsf{CC}}^\sharp\).  Hence
\(\mathcal H^\sharp,s^\sharp\entails\varphi^\sharp\), and by (c)
\(\mathcal H,s\entails\varphi\).  So
\(\Gamma\entails_{\mathrm H}^{\mathsf{CC}}\varphi\).

\emph{Backward.}  Assume
\(T_{\mathsf{CC}}^\sharp\cup\Gamma^\sharp\not\derives_{\mathsf{MSFOL}}\varphi^\sharp\),
that is,
\(\mathsf F\not\derives_{\mathsf{ND}_{\mathsf{FOL}}}\operatorname{tr}_{\mathrm{MS}}(\varphi^\sharp)\).
By \cref{cor:fol-completeness-countable} (the flat signature is countable),
\(\mathsf F\not\entails\operatorname{tr}_{\mathrm{MS}}(\varphi^\sharp)\), and by
\cref{thm:p02-many-sorted-translation},
\(T_{\mathsf{CC}}^\sharp\cup\Gamma^\sharp\not\entails_{\mathrm{MS}}\varphi^\sharp\):
there are a strict many-sorted structure \(\mathcal A\) and a sorted
assignment \(s\) with \(\mathcal A,s\entails\theta\) for every
\(\theta\in T_{\mathsf{CC}}^\sharp\), \(\mathcal A,s\entails\psi^\sharp\) for
every \(\psi\in\Gamma\), and \(\mathcal A,s\not\entails\varphi^\sharp\).  By
\cref{thm:p02-w02-p02-041-collapse}(a) and (d), \(\mathcal H_{\mathrm{ext}}\)
is an \(\mathsf H_{\mathsf{CC}}\)-model and \(s^{\mathrm{ext}}\) is matching;
by (c), \(\mathcal H_{\mathrm{ext}},s^{\mathrm{ext}}\entails\psi\) for every
\(\psi\in\Gamma\) and
\(\mathcal H_{\mathrm{ext}},s^{\mathrm{ext}}\not\entails\varphi\).  So
\(\Gamma\not\entails_{\mathrm H}^{\mathsf{CC}}\varphi\).  Contraposition gives
the backward direction.
\end{proof}
\end{annexitem}

\section{Proof for the formulas 30 of Church's system}
\label{sec:annex-p02-w02-p02-041-church}

\begin{annexitem}{annex:p02-w02-p02-041-church-30}{thm:p02-w02-p02-041-church-30}{Not every formula 30 is a theorem}
\begin{proof}
Work in \(\mathcal S(\N,c)\) with the denotation \(\sem{\cdot}_s\) of
\cref{def:p02-w02-p02-040-church-standard}.  For a type \(\beta\), write
\(\varphi^k\) for the \(k\)-fold composite of
\(\varphi\in\mathcal D_{\beta\beta}\) (\(\varphi^0\) the identity) and
\(\bar k^{\beta}\) for the member \(\varphi\mapsto\varphi^k\) of
\(\mathcal D_{\beta'}\), recalling that \(\beta'\) is
\(((\beta\beta)(\beta\beta))\).

\emph{Step 1: numerals, successor and \(N\).}  From the abbreviations,
\(k_{\beta'}\) is \(\lambda f_{\beta\beta}\lambda x_\beta\) followed by \(k\)
applications of \(f_{\beta\beta}\) to \(x_\beta\), so the denotation clauses
give \(\sem{k_{\beta'}}_s=\bar k^\beta\) for every \(s\).  Since
\(S_{\beta'\beta'}\) is
\(\lambda n_{\beta'}\lambda f_{\beta\beta}\lambda x_\beta(f_{\beta\beta}(n_{\beta'}f_{\beta\beta}x_\beta))\),
its denotation \(\sigma_\beta\) satisfies
\(\sigma_\beta(m)(\varphi)(v)=\varphi(m(\varphi)(v))\); hence
\(\sigma_\beta(\bar k^\beta)=\overline{k+1}^\beta\), and by induction on \(k\),
\(\sigma_\beta^k(\bar0^\beta)=\bar k^\beta\).  Unfolding the abbreviation of
\(N_{o\beta'}\) with the clauses for \(\Pi\), \(\sim\) and \(\vee\), and using
that \(\mathcal D_{o\beta'}\) contains the characteristic function of every
subset of \(\mathcal D_{\beta'}\), \(\sem{N_{o\beta'}}_s(m)=\mathrm T\) iff
\(m\) belongs to every subset of \(\mathcal D_{\beta'}\) that contains
\(\bar0^\beta\) and is closed under \(\sigma_\beta\).  Each such subset
contains \(\bar k^\beta\) for every \(k\), by induction on \(k\) using
\(\sigma_\beta(\bar k^\beta)=\overline{k+1}^\beta\); so
\(\sem{N_{o\beta'}}_s(\bar k^\beta)=\mathrm T\) for every \(k\).  By
\cref{thm:p02-w02-p02-040-church-soundness}(1),
\([\mathbf A_\gamma=\mathbf B_\gamma]\) denotes \(\mathrm T\) iff the two
denotations are equal.

\emph{Step 2: the formula \(m_{o''}S_{o'o'}0_{o'}\).}  Take \(\beta=o'\), so
that \(\beta'=o''\).  If \(s(m_{o''})=\bar k^{o'}\), then
\[
 \sem{m_{o''}S_{o'o'}0_{o'}}_s
 =\bar k^{o'}(\sigma_o)(\bar0^o)
 =\sigma_o^k(\bar0^o)
 =\bar k^{o},
\]
by Step 1 at \(\beta=o\).

\emph{Step 3: \(\bar1^o=\bar3^o\).}  \(\mathcal D_{oo}\) consists of the four
functions on \(\{\mathrm T,\mathrm F\}\): the identity, Boolean negation and
the two constant functions.  Each equals its own third iterate: the identity
and the constants are idempotent, and negation composed with itself is the
identity.  So \(\varphi^1=\varphi^3\) for every \(\varphi\in\mathcal D_{oo}\),
that is, \(\bar1^o=\bar3^o\) in \(\mathcal D_{o'}\).

\emph{Step 4: \(\bar1^{o'}\ne\bar3^{o'}\).}  \(\mathcal D_{o'}\) is the set of
functions from \(\mathcal D_{oo}\) to itself and has \(4^4=256\) members.
Choose distinct \(u_0,u_1,u_2\in\mathcal D_{o'}\) and let \(\psi\in\mathcal
D_{o'o'}\) send \(u_0\mapsto u_1\), \(u_1\mapsto u_2\), \(u_2\mapsto u_0\) and
fix every other member.  Then \(\psi^3\) is the identity while
\(\psi(u_0)=u_1\ne u_0\), so
\(\bar1^{o'}(\psi)=\psi\ne\psi^3=\bar3^{o'}(\psi)\), and
\(\bar1^{o'}\ne\bar3^{o'}\) in \(\mathcal D_{o''}\).

\emph{Step 5: a falsifying assignment.}  Let \(s\) satisfy
\(s(m_{o''})=\bar1^{o'}\) and \(s(n_{o''})=\bar3^{o'}\).  By Step 1,
\(N_{oo''}m_{o''}\) and \(N_{oo''}n_{o''}\) denote \(\mathrm T\).  By Steps 2
and 3 the two sides of
\(m_{o''}S_{o'o'}0_{o'}=n_{o''}S_{o'o'}0_{o'}\) denote \(\bar1^o=\bar3^o\), so
the equation denotes \(\mathrm T\).  By Step 4,
\(m_{o''}=n_{o''}\) denotes \(\mathrm F\).  By the dot convention 30\(^o\) is
\([N_{oo''}m_{o''}\supset[N_{oo''}n_{o''}\supset[\mathbf E\supset
m_{o''}=n_{o''}]]]\), with \(\mathbf E\) the equation, and
\([\mathbf A\supset\mathbf B]\) is \([[\sim\mathbf A]\vee\mathbf B]\), whose
denotation is \(\mathrm F\) exactly when \(\mathbf A\) denotes \(\mathrm T\)
and \(\mathbf B\) denotes \(\mathrm F\).  Working from the inside, the three
implications denote \(\mathrm F\).  So 30\(^o\) is not valid in
\(\mathcal S(\N,c)\).

\emph{Conclusion.}  By \cref{thm:p02-w02-p02-040-church-soundness}(2) every
theorem is valid in \(\mathcal S(\N,c)\); hence 30\(^o\) is not a theorem,
and not all formulas 30\(^\alpha\) are theorems.
\end{proof}
\end{annexitem}
